Let \(\widetilde T_n\) denote the expression inside the clipping operation and let \(\mathcal G:=\sigma((X_i,A_i)_{i=1}^n)\).  Lemma~\(\mathrm{lem:scale\text{-}sanity}\) gives \(\tau(P)\in[-M,M]\), so projection onto this interval cannot increase squared loss.  It is enough to control \(\widetilde T_n\).

Conditional independence and the definitions of the cell means give
\[
 E[\widetilde T_n\mid\mathcal G]
 =\frac{\mathbf 1\{R_n>0\}}{R_n}\sum_{k=1}^dN_kI_k\tau_k.
\tag{1}
\]
On \(\{R_n>0\}\), subtracting \(\tau(P)\) from \((1)\) produces a convex combination of the deviations \(\delta_k\), so approximate homogeneity bounds its absolute value by \(\sigma M\).  On \(\{R_n=0\}\), the conditional mean of the fallback is zero and \(|\tau(P)|\le M\).  Hence
\[
 E\!\left[\{E[\widetilde T_n\mid\mathcal G]-\tau(P)\}^2\right]
 \le M^2\{\sigma^2+Q_P^{(n)}(R_n=0)\}.
\tag{2}
\]
Write \(v_{ak}:=\operatorname{Var}(Y\mid A=a,X=k)\).  Assumption~\(\mathrm{ass:second\text{-}central\text{-}moment}\) gives \(v_{ak}\le M^2\).  Conditional on \(\mathcal G\), the empirical arm means use disjoint observations and therefore
\[
 \operatorname{Var}(\widetilde T_n\mid\mathcal G)
 \le \frac{M^2\mathbf 1\{R_n>0\}}{R_n^2}
 \sum_{k=1}^dN_k^2I_k\left(\frac1{N_{1k}}+\frac1{N_{0k}}\right).
\tag{3}
\]
Every denominator in \((3)\) is positive because its summand is multiplied by \(I_k\).

Conditional on the category counts, the variables \(N_{1k}\) are independent \(\operatorname{Binomial}(N_k,\pi_k)\) variables.  If \(D\sim\operatorname{Binomial}(m,q)\), \(m\ge2\), and \(J:=\mathbf 1\{0<D<m\}\), then
\[
 E\left[\frac1D\,\middle|\,J=1\right]
 \le \frac{2}{(m+1)q\{1-q^m-(1-q)^m\}}
 \le \frac{C_{\epsilon}}m,
\tag{4}
\]
uniformly for \(q\in[\epsilon,1-\epsilon]\).  Indeed, \(D^{-1}\le2(D+1)^{-1}\) on \(D\ge1\),
\[
 E[(D+1)^{-1}]=\frac{1-(1-q)^{m+1}}{(m+1)q},
\]
and \(P(J=1)=1-q^m-(1-q)^m\ge2q(1-q)\ge2\epsilon(1-\epsilon)\).  The same calculation applies to \((m-D)^{-1}\).  Conditioning further on \((I_k)_k\) preserves independence across cells.  Since \(R_n=\sum_kN_kI_k\), equations~\((3)\)--\((4)\) yield
\[
 E[\operatorname{Var}(\widetilde T_n\mid\mathcal G)\mid(N_k,I_k)_k]
 \le C_{\epsilon}M^2\frac{\mathbf 1\{R_n>0\}}{R_n}.
\tag{5}
\]
The conditional bias--variance identity, \((2)\), \((5)\), and Lemma~\(\mathrm{lem:zeng\text{-}usable\text{-}occupancy\text{-}reciprocal}\) imply
\[
 E[(\widetilde T_n-\tau(P))^2]
 \le C_{\epsilon}M^2\left\{\sigma^2+\frac{n\vee d}{n^2}
 +\exp\left(-b_{\epsilon}\frac{n^2}{n\vee d}\right)\right\}.
\tag{6}
\]
If \(d\le n\), the exponential term is at most \(C_{\epsilon}/n\).  If \(n<d\le c_{\epsilon}n^2/\log(en)\), decreasing \(c_{\epsilon}\), if necessary, makes that term at most \(C_{\epsilon}d/n^2\).  Finally, \((n\vee d)/n^2\le n^{-1}+d/n^2\).  Substitution in \((6)\), followed by clipping, proves the displayed bound.  Conditional moments are invoked only for positive-mass cells, and all empirical division boundaries have the stipulated zero fallback.